\section{Discussion and Limitations}
\subsection{When compensation-aware ranking is useful}
The mechanism is expected to be most relevant when compensation workspace or response bandwidth differs across gait candidates. In an unconstrained regime with a near-perfect stabilizer, the extra gait selection criterion may produce little benefit. That regime is a legitimate limiting case and should be included in analysis rather than hidden.

The model can also fail when predicted disturbances are dominated by uncertain contact events or when the optimal gait lies outside the small candidate set. A physically reliable response must explicitly address low-bandwidth position servos, contact uncertainty and the limited directions of body correction. Merely allowing a six-dimensional virtual camera correction would overstate what the legged platform can accomplish.

\subsection{Measurement limitations and real-world interpretation}
Camera twist prediction can be related to image displacement only through a stated camera model and visible depth distribution. Inter-frame jitter, exposure-time path length, motion blur, rolling-shutter shear and perception accuracy are not interchangeable metrics. A feature-depth weighting that improves short-range targets may not minimize background image motion; this tradeoff should be discussed for the eventual task. The real-camera comparison is necessary for any explicit claim about recorded video stabilization.

\subsection{Scope of claimed contribution}
The work cannot claim to be the first hexapod controller with gaze feedback, the first gait smoother, the first CPG--MPC hierarchy, or the first visual objective in robot control. Any eventual novelty claim must be supported by measured improvement over residual-image-cost scheduling with a candidate-independent compensation budget, not solely by superiority to a naive fixed tripod gait.
